\begin{enumerate}[label=\textbf{\ansnum.\arabic*}]
  \item An interpreter reads the program text and executes it while the program runs; a compiler translates the whole program once into machine code before it runs, and the result runs by itself. For a compiled Rust program nothing has to be installed on the robot except the executable file (and a Python program needs Python and its packages).
  \item $2\,\mathrm{ms}/229\,\mathrm{ns}\approx8700$ steps in Python; $2\,\mathrm{ms}/8.9\,\mathrm{ns}\approx225\,000$ steps in Rust, which is $26$ times more.
  \item Rust: $3\cdot10^9\cdot8.9\,\mathrm{ns}=26.7$\,s. Python: $3\cdot10^9\cdot229\,\mathrm{ns}=687$\,s, about $11.5$ minutes.
  \item $1828/217\approx8.4$ times faster. The debug build compiles much faster and contains information for a debugger and extra checks (for example for numeric overflow), which is what you want while writing and testing; the release build is for measurements and for the robot.
  \item Use after free (a pointer to the data of a camera frame that was already released: wrong ball position or a crash); buffer overflow (writing the 23rd value into an array of $22$ joint values: neighbouring data is overwritten, a joint receives a wrong command); data race (the vision thread writes the ball position while the control thread reads it: half-new, half-old coordinates); uninitialised variable (a gain that was never set is $0$ or random: the robot does not react or moves wildly). Other sensible examples are accepted.
  \item $30\,\mathrm{ms}\cdot200\,\mathrm{Hz}=6$ steps (the period is $5$\,ms). During six periods nothing is measured or commanded: a balance controller can let the robot fall, a kick can be missed.
  \item (a) Python: plots, notebooks, quick changes. (b) Rust: fixed deadline, no pauses, runs for a whole match without crashes. (c) Python (or a shell script): written once, speed irrelevant. (d) Rust: delivered as a single file to many laptops, the same language and data types as the robot program, fast enough to show data at high rate in the terminal.
  \item A wrong formula or logic error, for example using degrees where the function expects radians, swapping left and right, or a sign error in a controller: it compiles and runs. One needs tests with known results (Chapter~\ref{ch:rust5}), simulation, and checks with the real robot.
  \item Cross-compilation is compiling on one kind of computer (the \emph{host}, your laptop) a program for another kind (the \emph{target}, the robot). The build is fast because it runs on the faster machine, you need no compiler or source code on the robot, and you can build the program in the same way for each student. The details are in Chapter~\ref{ch:rust10}.
  \item[\textbf{\ansnum.10}$\star$] For: speed of loops, one executable, safer long runs. Against: the project is dominated by libraries that are already compiled code (OpenCV, NumPy, video decoding) so Rust would not make it much faster; analysis needs plots, notebooks and quick experiments, which are Python's strength; the rewrite takes weeks and the team knows Python best. Recommendation: keep Python, and rewrite only a part if a measurement shows that this part is the bottleneck. First measurement: a profile of the existing program (for example with \verb!cProfile!, Chapter~\ref{ch:py11}) that shows where the time is spent, and whether the program is waiting for the disk or the video decoder rather than computing.
\end{enumerate}
